\subsection*{Partie A}

On considère l'équation différentielle 

\[(E_1) :\quad  y' + 0{,}48y = \dfrac{1}{250},\]

où $y$ est une fonction de la variable $t$ appartenant à l'intervalle $[0~;~ +\infty[$.

\begin{enumerate}
\item On considère la fonction constante $h$ définie sur l'intervalle $[0~;~ +\infty[$ par $h(t) = \dfrac{1}{120}$.

Montrer que la fonction $h$ est solution de l'équation différentielle $(E_1)$.
\item Donner la forme générale des solutions de l'équation différentielle $y' + 0,48y = 0$.
\item En déduire l'ensemble des solutions de l'équation différentielle $(E_1)$.
\end{enumerate}

\subsection*{Partie B}

On s'intéresse à présent à l'évolution d'une population de bactéries dans un milieu de culture.

À un instant $t = 0$, on introduit une population initiale de \np{30000} bactéries dans le milieu. On note $p(t)$ la quantité de bactéries, exprimée en millier d'individus, présente dans le milieu après un temps $t$, exprimé en heure.

On a donc $p(0) = 30$.

On admet que la fonction $p$ définie sur l'intervalle $[0~;~ +\infty[$ est dérivable, strictement positive sur cet intervalle et qu'elle est solution de l'équation différentielle $(E_2)$ :

\[p' = \dfrac{1}{250}p \times (120 - p)\]

Soit $y$ la fonction strictement positive sur l'intervalle $[0~;~ +\infty[$ telle que, pour tout $t$ appartenant à l'intervalle $[0~;~ +\infty[$ , on a $p(t) = \dfrac{1}{y(t)}$.

\medskip

\begin{enumerate}
\item Montrer que si $p$ est solution de l'équation différentielle $(E_2)$, alors $y$ est solution de l'équation différentielle $(E_1)$ : \quad $y' + 0,48y = \dfrac{1}{250}$.
\item On admet réciproquement que, si $y$ est une solution strictement positive de l'équation différentielle $(E_1)$, alors $p = \dfrac 1y$ est solution de l'équation différentielle $(E_2)$.

Montrer que, pour tout $t$ appartenant à l'intervalle $[0~;~ +\infty[$ , on a :

\[p(t) = \dfrac{120}{1 + K\e^{-0,48t}}~ \text{avec}\: K\: \text{une constante réelle}.\] 
\item En utilisant la condition initiale, déterminer la valeur de $K$.
\item Déterminer $\displaystyle\lim_{t \to + \infty} p(t)$. En donner une interprétation dans le contexte de l'exercice.
\item Déterminer le temps nécessaire pour que la population de bactéries dépasse \np{60000}~individus.

\emph{On donnera le résultat sous la forme d'une valeur arrondie exprimée en heures et minutes.}
\end{enumerate}
